\documentclass[10pt,landscape]{article}
\usepackage{amssymb,amsmath,amsthm,amsfonts}
\usepackage{multicol,multirow}
\usepackage{marvosym}
\usepackage{calc}
\usepackage{ifthen}
\usepackage[landscape]{geometry}
\usepackage[colorlinks=true,citecolor=blue,linkcolor=blue]{hyperref}
\usepackage{notes_2023}

\setlength{\extrarowheight}{0pt}

\ifthenelse{\lengthtest { \paperwidth = 11in}}
{ \geometry{top=.5in,left=.5in,right=.5in,bottom=.5in} }
{\ifthenelse{ \lengthtest{ \paperwidth = 297mm}}
	{\geometry{top=1cm,left=1cm,right=1cm,bottom=1cm} }
	{\geometry{top=1cm,left=1cm,right=1cm,bottom=1cm} }
}
%\pagestyle{empty}
\makeatletter
\renewcommand{\section}{\@startsection{section}{1}{0mm}%
	{-1ex plus -.5ex minus -.2ex}%
	{0.5ex plus .2ex}%x
	{\normalfont\large\bfseries}}
\renewcommand{\subsection}{\@startsection{subsection}{2}{0mm}%
	{-1explus -.5ex minus -.2ex}%
	{0.5ex plus .2ex}%
	{\normalfont\normalsize\bfseries}}
\renewcommand{\subsubsection}{\@startsection{subsubsection}{3}{0mm}%
	{-1ex plus -.5ex minus -.2ex}%
	{1ex plus .2ex}%
	{\normalfont\small\bfseries}}
\makeatother
\setcounter{secnumdepth}{0}
\setlength{\parindent}{0pt}
\setlength{\parskip}{0pt plus 0.5ex}
% -----------------------------------------------------------------------

\title{Scheda riassuntiva di Teoria dei campi e di Galois}

\begin{document}

\parskip=0.7ex

\raggedright
\footnotesize

\begin{center}
	\Large{\textbf{Scheda riassuntiva di Teoria dei campi e di Galois}} \\
\end{center}
\begin{multicols}{3}
	\setlength{\premulticols}{1pt}
	\setlength{\postmulticols}{1pt}
	\setlength{\multicolsep}{1pt}
	\setlength{\columnsep}{2pt}

	\textit{La teoria si concentra sulle estensioni finite su campi perfetti,
		benché si sia dato spazio anche a considerazioni su estensioni infinite
		o non separabili.}

	\section{Campi e omomorfismi}

	Si dice \textbf{campo} un anello commutativo non banale
	$K$ i cui elementi non nulli ammettono un inverso
	moltiplicativo. Si dice
	\textbf{omomorfismo di campo} tra due campi $K$ ed $L$
	un omomorfismo che è anche omomorfismo di anelli. \smallskip

	Dal momento che un omomorfismo
	$\varphi$ è tale per cui $\Ker \varphi$ è un ideale
	di $K$ con $1 \notin \Ker \varphi$, deve per forza
	valere $\Ker \varphi = \{0\}$, e quindi ogni omomorfismo
	di campi è sempre iniettivo.

	\section{Caratteristica di un campo}

	Esiste un unico omomorfismo $\ZZ \to K$ completamente
	determinato dalla relazione $1 \mapsto 1_K$, e
	si definisce \textbf{caratteristica di $K$}, detta
	$\Char K$, il
	generatore non negativo del nucleo di questa mappa. In particolare
	$\Char K$ è $0$ o un numero primo. \smallskip

	Se $\Char K$ è zero,
	$\zeta$ è un'immersione, e quindi $K$ è un campo infinito,
	e in particolare vi si immerge anche $\QQ$. \smallskip

	Tuttavia non è detto che $\Char K = p$ implichi che $K$ è
	finito. In particolare $\FF_p(x)$, il campo delle funzioni
	razionali a coefficienti in $\FF_p$, è un campo infinito
	a caratteristica $p$.

	\subsection{Proprietà dei campi a caratteristica \texorpdfstring{$p$}{p}}

	Se $\Char K = p$, per il Primo
	teorema di isomorfismo per anelli, $\FF_p$ si immerge
	su $K$ tramite il passaggio al quoziente di $\ZZ \to K$; pertanto
	$K$ contiene una copia isomorfa di $\FF_p$. Per
	campi di caratteristica $p$, vale il Teorema del
	binomio ingenuo, ossia:
	\[ (a + b)^p = a^p + b^p, \]
	estendibile anche a più addendi. \smallskip

	In particolare, per un campo $K$ di caratteristica $p$,
	la mappa $\Frob : K \to K$ tale per cui $a \xmapsto{\Frob} a^p$
	è un omomorfismo di campi, ed in particolare è un'immersione
	di $K$ in $K$, detta \textbf{endomorfismo di Frobenius}. \smallskip

	Se $K$ è un campo finito, $\Frob$ è anche
	un isomorfismo (gli omomorfismi tra campi sono iniettivi!).
	Si osserva che per gli elementi di $\FF_p \subseteq K$ vale
	$\restr{\Frob}{\FF_p} = \Id_{\FF_p}$, e quindi
	$\Frob$ è un elemento di $\Gal(K / \FF_p)$.

	\section{Campi finiti}

	Per ogni $p$ primo e $n \in \NN^+$ esiste un campo finito
	di ordine $p^n$. Ogni tale campo può essere visto
	come spazi vettoriale di dimensione $n$ sull'immersione di $\FF_p$ che contiene,
	e come campo di spezzamento di $x^{p^n}-x$
	su tale immersione. Tale campi ha obbligatoriamente
	caratteristica $p$. \smallskip

	Fissati $p$ ed $n$, esiste un unico campo finito di cardinalità
	$p^n$: esiste sempre un isomorfismo tra due campi finiti che
	fissa le copie di $\FF_p$. Indichiamo con $\FF_{p^n}$ la struttura
	algebrica di un campo di cardinalità $p^n$. \smallskip

	Vale la relazione $\FF_{p^n} \subseteq \FF_{q^m}$
	se e solo se $p=q$ e $n \mid m$ (segue facilmente dal Teorema
	delle torri algebriche). Conseguentemente,
	l'estensione minimale per inclusione comune a
	$\FF_{p^{n_1}}$, ..., $\FF_{p^{n_i}}$ è
	$\FF_{p^m}$ dove $m := \mcm(n_1, \ldots, n_i)$. Pertanto
	se $p \in \FF_{p^n}[x]$ si decompone in fattori irriducibili
	di grado $n_1$, \ldots, $n_i$, il suo campo di spezzamento
	è $\FF_{p^m}$. Da ciò segue anche che $x^{p^n}-x$ è in $\FF_p$ il
	prodotto di tutti gli irriducibili di grado divisore
	di $n$.

	\section{Campi di spezzamento per polinomi su \texorpdfstring{$K$}{K}}

	Per il Teorema di Lagrange sui campi, ogni polinomio
	di $K[x]$ ammette al più tante radici quante il suo grado.
	Come conseguenza pratica di questo teorema, ogni sottogruppo
	moltiplicativo finito di $K$ è ciclico. Pertanto
	$\FF_{p^n}^* = \gen{\alpha}$ per $\alpha \in \FF_{p^n}$,
	e quindi $\FF_{p^n} = \FF_p(\alpha)$, ossia
	$\FF_{p^n}$ è sempre un'estensione semplice su $\FF_p$. \smallskip

	Si dice
	\textbf{campo di spezzamento} di una famiglia $\mathcal{F}$
	di polinomi di $K[x]$ un sovracampo minimale per
	inclusione di $K$ che fa sì che ogni polinomio di $\mathcal{F}$
	si decomponga in fattori lineari. I campi
	di spezzamento di $\mathcal{F}$ sono sempre
	$K$-isomorfi tra loro (ossia ammettono isomorfismi che fissano
	il campo base $K$). \smallskip

	Un polinomio irriducibile si dice separabile se ammette
	radici distinte. Per il Criterio della derivata,
	$p \in K[x]$ ammette radici multiple se e solo se
	$\MCD(p, p')$ non è invertibile, dove $p'$ è la derivata
	formale di $p$. \smallskip

	Se $p \in K[x]$ e $n := \deg p$, il campo di spezzamento $L$ di $p$ è tale per cui
	$[L : K] \leq n!$ (infatti, presa una radice, la sua aggiunta al campo base
	restituisce un'estensione di grado al più
	$n$; dobbiamo poi considerare il polinomio rimanente di grado $n-1$ che aumenta l'estensione
	di al più un fattore $n-1$; reiterando si ottiene la tesi). Se $p$ è irriducibile e separabile, vale anche che $n \mid [L : K] \mid n!$, come
	conseguenza dell'azione del relativo gruppo di
	Galois sulle radici. \smallskip

	Se $p$ è irriducibile in $K[x]$, $(p)$ è un ideale
	massimale, e $K[x] / (p)$ è un campo che
	ne contiene una radice, ossia $\overline x$. In
	particolare $K$ si immerge in $K[x] / (p)$,
	e quindi tale campo può essere identificato come
	un'estensione di $K$ che aggiunge una radice di $p$.
	Se $K$ è finito, detta $\alpha$ la radice aggiunta
	all'estensione, $L := K[x] / (p) \cong K(\alpha)$ contiene
	tutte le radici di $p$ (ed è dunque il suo campo
	di spezzamento). Infatti detto $[L : \FF_p] = n$,
	$\overline x$ annulla $x^{p^n}-x$ per il Teorema di Lagrange
	sui gruppi, e quindi $p$ deve dividere $x^{p^n}-x$;
	in tal modo $p$ deve spezzarsi in fattori lineari,
	e quindi ogni radice deve già appartenere ad $L$
	(alternativamente, esiste un unico campo di cardinalità $p^n$ a
	meno di isomorfismo, e quindi necessariamente le radici devono
	già appartenere in $\FF_{p^n}$).
	In particolare, ogni estensione finita e semplice
	di un campo finito è normale, e quindi di Galois.

	\section{Estensioni di campo}

	Si dice che $L$ è un'estensione di $K$, e si indica
	con $L / K$, se $L$ è un sovracampo di $K$,
	ossia se $K \subseteq L$. Si indica con $[L : K] = \dim_K L$ la
	dimensione di $L$ come $K$-spazio vettoriale. Si
	dice che $L$ è un'estensione finita di $K$ se $[L : K]$
	è finito, e infinita altrimenti. Un'\textbf{estensione finita}
	di un campo finito è ancora un campo finito. Un'estensione
	è finita se e solo se è finitamente generata da elementi algebrici. \smallskip

	Una $K$-immersione è un omomorfismo di campi
	da un'estensione di $K$ in un'altra estensione di $K$ che
	agisce come l'identità su $K$. Un $K$-isomorfismo è
	una $K$-immersione che è isomorfismo.

	\subsection{Composto di estensioni e teorema delle torri algebriche}

	Date estensioni $L$ e $M$ su $K$, si definisce
	$LM = L(M) = M(L)$ come il \textbf{composto} di $L$
	ed $M$, ossia come la più piccola estensione di $K$ che
	contiene sia $L$ che $M$. In particolare, $LM$
	può essere visto come $L$-spazio vettoriale con
	vettori in $M$, o analogamente come $M$-spazio con
	vettori in $L$. \smallskip

	Per il Teorema delle torri algebriche, $L / K$ è
	un'estensione finita se e solo se $L / F$ e
	$F / K$ lo sono (ossia la finitezza vale strettamente
	per torri). Inoltre, se $\basis_{L/F}$ e $\basis_{F/K}$
	sono basi di $L/F$ e $F/K$, allora
	$\basis_{L/F} \basis_{F/K}$ è una base di
	$L / K$, dove i suoi elementi sono i prodotti tra
	i vari elementi delle due basi. Infine
	se $L / K$ è finita, allora
	anche $LM / M$ è finita, e vale che $[LM : M] \leq [L : K]$ (infatti una base di $L / K$ può essere trasformata
	in un insieme di generatori di $LM / M$), e quindi
	la finitezza vale per \textit{shift}. Sempre per
	il Teorema delle torri algebriche, se $L / K$ è
	finito, allora vale che:
	\[ [L : K] = [L : F] [F : K]. \]
	Se $L / K$ e $M / K$ sono finite, anche $LM / K$ lo
	è (infatti la finitezza vale sia per torri che per \textit{shift}). In particolare, vale che:
	\[ \mcm([L : K], [M : K]) \mid [LM : K]. \]
	Se $[L : K]$ ed $[M : K]$ sono coprimi tra loro,
	allora vale proprio l'uguaglianza
	$[LM : K] = [L : K] [M : K]$. Infatti, in tal caso,
	si avrebbe $[L : K] [M : K] \leq [LM : K]$ e
	$[LM : K] = [LM : M] [M : K] \leq [L : K] [M : K]$.

	\subsection{Omomorfismo di valutazioni, elementi algebrici e trascendenti e polinomio minimo}

	Dato $\alpha$, si definisce $K(\alpha)$ il più piccolo
	sovracampo di $K$ che contiene $\alpha$. Si definisce l'\textbf{omomorfismo di
		valutazione} $\varphi_{\alpha, K} : K[x] \to K[\alpha]$, detto
	$\varphi_\alpha$ se $K$ è noto, l'omomorfismo completamente
	determinato dalla relazione $p \xmapsto{\varphi_\alpha} p(\alpha)$. Si verifica che $\varphi_\alpha$ è
	surgettivo. Se $\varphi_\alpha$ è iniettivo,
	si dice che $\alpha$ è \textbf{trascendente} su $K$ e
	$K[x] \cong K[\alpha]$, da cui $[K[\alpha] : K] =
		[K[x] : K] = \infty$. Se invece $\varphi_\alpha$ non
	è iniettivo, si dice che $\alpha$ è \textbf{algebrico}
	su $K$. Si definisce $\mu_\alpha$, detto il \textbf{polinomio
		minimo} di $\alpha$ su $K$, il generatore monico
	di $\Ker \varphi_\alpha$. Dal momento che $K$ è
	in particolare un dominio di integrità e che $K[x]$ è un PID,
	$\mu_\alpha$ è sempre irriducibile. \smallskip

	Si definisce
	$\deg_K \alpha := \deg \mu_\alpha$. Se $\alpha$ è
	algebrico su $K$, $K[x] / (\mu_\alpha) \cong
		K[\alpha]$, e quindi $K[\alpha]$ è un campo. Dacché
	$K[\alpha] \subseteq K(\alpha)$, vale allora
	$K[\alpha] = K(\alpha)$. Inoltre, poiché $\dim_K K[x] / (\mu_\alpha) = \deg_K \alpha$, vale
	anche che $[K(\alpha) : K] = \deg_K \alpha$.
	Infine, si verifica che $\alpha$ è algebrico se e solo se
	$[K(\alpha) : K]$ è finito. \smallskip

	\subsection{Estensioni semplici, algebriche}

	Si dice che $L$ è un'\textbf{estensione semplice} di
	$K$ se $\exists \alpha \in L$ tale per cui $L = K(\alpha)$.
	In tal caso si dice che $\alpha$ è un \textbf{elemento primitivo} di $K$. Si dice che $L$ è un'\textbf{estensione
		algebrica} di $K$ se ogni suo elemento è algebrico su $K$.
	Ogni estensione finita è algebrica. Non tutte le
	estensioni algebriche sono finite (e.g.~$\overline{\QQ}$ su $\QQ$). \smallskip

	Gli elementi algebrici di $K$ formano un'estensione algebrica di $K$.
	Pertanto se $\alpha$ e $\beta$ sono algebrici,
	$\alpha \pm \beta$, $\alpha \beta$, $\alpha \beta\inv$
	e $\alpha\inv \beta$ (a patto che opportunamente $\alpha \neq 0$ o
	$\beta \neq 0$) sono algebrici. \smallskip

	Dati $a$, $b \in \QQ$, le estensioni $\QQ(\sqrt{a})$ e
	$\QQ(\sqrt{b})$ sono uguali se e solo se $ab$ (o equivalentemente
	$a/b$) è un quadrato in $\QQ$.

	\subsection{Campi perfetti, estensioni separabili e coniugati}

	Si dice che un'estensione algebrica $L$ è un'\textbf{estensione separabile} di
	$K$ se per ogni elemento $\alpha \in L$,
	$\mu_\alpha$ ammette radici distinte. Si dice
	che $K$ è un \textbf{campo perfetto} se ogni
	polinomio irriducibile ammette radici distinte,
	ossia se ogni polinomio irriducibile è separabile.
	In un campo perfetto, ogni estensione algebrica
	è separabile. Si definiscono i \textbf{coniugati} di
	$\alpha$ algebrico su $K$ come le radici
	di $\mu_\alpha$. Se $K(\alpha)$ è separabile su $K$,
	$\alpha$ ha esattamente $\deg_K \alpha$ coniugati,
	altrimenti esistono al più $\deg_K \alpha$ coniugati. \smallskip

	Un campo è perfetto se e solo se ha caratteristica
	$0$ o altrimenti se l'endomorfismo di
	Frobenius è un automorfismo. Equivalentemente,
	un campo è perfetto se le derivate dei polinomi
	irriducibili sono sempre non nulle. Esempi di
	campi perfetti sono allora tutti i campi di
	caratteristica $0$ e tutti i campi finiti.

	\subsection{Campi algebricamente chiusi e chiusura algebrica di \texorpdfstring{$K$}{K}}

	Un campo $K$ si dice \textbf{algebricamente chiuso} se
	ogni $p \in K[x]$ ammette una radice in $K$. Equivalentemente $K$ è algebricamente chiuso se
	ogni $p \in K[x]$ ammette tutte le sue radici in $K$.
	Si dice \textbf{chiusura algebrica} di $K$
	una sua estensione algebrica e algebricamente
	chiusa. Le chiusure algebriche di $K$ sono
	$K$-isomorfe tra loro, e quindi si identifica
	con $\overline{K}$ la struttura algebrica della
	chiusura algebrica di $K$. \smallskip

	Se $L$ è una sottoestensione di $K$ algebricamente
	chiuso, allora $\overline{L}$ è il campo degli
	elementi algebrici di $K$ su $L$. Infatti se
	$p \in L[x]$, $p$ ammette una radice $\alpha$ in $K$, essendo
	algebricamente chiuso. Allora $\alpha$ è un elemento
	di $K$ algebrico su $L$, e quindi $\alpha \in \overline{L}$. Per il Teorema fondamentale dell'algebra,
	$\overline{\RR} = \CC$.

	\subsection{Estensioni normali e di Galois, \texorpdfstring{$K$}{K}-immersioni di un'estensione finita di \texorpdfstring{$K$}{K}}

	Sia $\alpha$ un elemento algebrico su $K$. Allora
	$[K(\alpha) : K] = \deg_K \alpha$. Le
	$K$-immersioni da $K(\alpha)$ in $\overline{K}$
	sono esattamente tante quanti sono i coniugati di
	$\alpha$ e sono tali da mappare $\alpha$ ad un suo coniugato. Infatti,
	ogni $K$-immersione si ottiene come passaggio al quoziente
	di una mappa $K[x] \to K$ che fissa $K$ il cui nucleo contiene
	$\mu_\alpha$. Se $K$ è perfetto, esistono esattamente
	$\deg_K \alpha$ $K$-immersioni da $K(\alpha)$
	in $\overline{K}$. \smallskip

	Generalizzando, se $L$ è un'estensione separabile finita, allora per
	ogni $\varphi : K \to \overline{K}$ esistono
	esattamente $[L : K]$ estensioni $\varphi_i : L \to \overline{K}$ di $\varphi$, ossia omomorfismi
	tali per cui $\restr{\varphi_i}{K} = \varphi$. \smallskip

	Se $L / K$ è un'estensione separabile finita su $K$, allora
	esistono esattamente $[L : K]$ $K$-immersioni
	da $L$ in $\overline{K}$. Per quanto detto prima,
	tali immersioni mappano gli elementi $L$ nei
	loro coniugati. \smallskip

	Per calcolare dunque tutti
	i coniugati di $\alpha \in L$ su $K$, è sufficiente
	calcolare i distinti valori delle $K$-immersioni
	di $L$ su $\alpha$. Infatti, ogni $K$-immersione
	da $K(\alpha)$ può estendersi a $K$-immersione di
	$L$, e viceversa ogni $K$-immersione di $L$ può
	restringersi a $K$-immersione di $K(\alpha)$. In
	particolare, una volta computati tutti i coniugati, è semplice trovare il polinomio minimo di $\alpha$
	su $K$ (è sufficiente considerare il prodotto dei vari $x-\alpha_i$ dove gli $\alpha_i$ sono tutti i coniugati di $\alpha$). \smallskip

	Si dice che un'estensione algebrica $L / K$ è un'\textbf{estensione normale}
	se per ogni $K$-immersione $\varphi$ da $L$ in $\overline{K}$
	vale che $\varphi(L) = L$. Equivalentemente
	un'estensione è normale se è il campo di spezzamento
	di una famiglia di polinomi (in particolare è il campo
	di spezzamento di tutti i polinomi irriducibili che
	hanno una radice in $L$). Ancora, un'estensione $L$
	è normale se e solo se per ogni $\alpha \in L$,
	i coniugati di $L$ appartengono ancora ad $L$. \smallskip

	Per un'estensione normale ad ogni $K$-immersione
	$\varphi : L \to \overline{K}$ si può restringere
	il codominio ad $L \subseteq \overline{K}$, e quindi considerare $\varphi$ come
	un automorfismo di $L$ che fissa $K$. \smallskip

	Se $\Char K \neq 2$, un'estensione finita $L/K$ di grado $2$ è sempre normale,
	ed in particolare può sempre scriversi come
	$L = K(\sqrt{\Delta})$, dove $\Delta$ non è un quadrato
	in $K$. \smallskip

	Si indica con $\Aut_K(L) = \Aut(L / K)$ l'insieme
	degli automorfismi di $L$ che fissano $K$. Se
	$L$ è normale e separabile, si dice che $L$ è
	un'\textbf{estensione di Galois}, e si definisce
	il suo \textbf{gruppo di Galois}
	$\Gal(L / K)$ come $(\Aut_K L, \circ)$, ossia come
	il gruppo $\Aut_K L$ con l'operazione di
	composizione. \smallskip

	\subsection{Azione di \texorpdfstring{$\Gal(L / K)$}{Gal(L/K)} sulle radici di \texorpdfstring{$L$}{L} campo di spezzamento}

	Sia $p \in K[x]$ irriducibile e separabile.
	Allora si definisce
	il \textbf{gruppo di Galois di $p$} come il gruppo
	di Galois $\Gal(L / K)$, dove $L$ è campo di
	spezzamento di $p$ su $K$. Se $\deg p = n$ e
	$a_1$, ..., $a_n$ sono le radici di $p$,
	$\Gal(L / K)$ agisce fedelmente e transitivamente
	su $\{a_1, \ldots, a_n\}$ in modo naturale:

	\begin{equation*}
		\begin{split}
			\Gal( & L / K) \curvearrowright S(\{a_1, \ldots, a_n\}) \cong S_n, \\
			      & \varphi_i \cdot a_j = \varphi_i(a_j).
		\end{split}
	\end{equation*}

	Poiché l'azione è fedele, vale che
	$\Gal(L / K) \mono S_n$, e quindi $\abs{\Gal(L / K)} \mid n!$, perché
	$n!$ è la cardinalità di $S_n$.
	Se $\Gal(L / K)$ è abeliano
	(e in tal caso si dice che $L$ è un'\textbf{estensione abeliana}), l'azione
	è anche libera, e quindi per il Lemma orbita-stabilizzatore vale
	$\abs{\Gal(L / K)} = n$. \smallskip

	Inoltre, poiché $[K(a_1) : K] = n$, vale in particolare:
	\[ n \mid \abs{\Gal(L / K)} = [L : K] \mid n!. \]

	\section{Proprietà su torri, traslato, composto e intersezione}

	Sia $\mathcal{P}$ una proprietà. Allora si
	studia la proprietà $\mathcal{P}$ secondo
	le seguenti quattro modalità:
	\begin{itemize}
		\item validità per \textbf{torri}: se $\mathcal{P}$ vale in due estensioni in $K \subseteq F \subseteq L$, allora vale anche per la terza estensione, ossia
		      vale per tutta la torre di estensioni,
		\item validità per \textbf{\textit{shift}} (o per il \textbf{traslato}): se $\mathcal{P}$ vale
		      per $L / K$, allora vale anche per $LM / M$, ossia
		      vale sul ramo parallelo a quello di $L / K$,
		\item validità per il \textbf{composto}: se
		      $\mathcal{P}$ vale per $L / K$ ed $M / K$, allora
		      vale anche per $LM / K$.
		\item validità per l'\textbf{intersezione}:
		      se $\mathcal{P}$ vale per $L / K$ ed $M / K$,
		      allora vale anche per $(L \cap M) / K$.
	\end{itemize}
	Si dice che $\mathcal{P}$ vale \textit{debolmente}
	per torri, se $\mathcal{P}$ vale per $L / K$ solo
	se vale per $L / F$ sottoestensione.
	Si dice che $\mathcal{P}$ vale \textit{strettamente}
	per torri, se è $\mathcal{P}$ vale per $L / K$ se
	e solo se vale per $L / F$ e $F / K$. Se $\mathcal{P}$ vale strettamente per torri, allora $\mathcal{P}$
	vale anche per l'intersezione. \smallskip

	Si dice che
	$\mathcal{P}$ vale \textit{inversamente} per
	\textit{shift} se $\mathcal{P}$ vale su
	$LF / F$ solo se vale su $L / K$. Si dice che
	$\mathcal{P}$ vale \textit{inversamente} per
	il composto se $\mathcal{P}$ vale su $LF / K$
	implica che $\mathcal{P}$ valga anche su $L / K$
	e $F / K$. Si dice che $\mathcal{P}$ vale \textit{completamente} per \textit{shift} o composto se $\mathcal{P}$
	vale \textit{inversamente} e normalmente per \textit{shift} o
	composto. Se $\mathcal{P}$ vale per torri e
	per \textit{shift}, allora vale anche per il
	composto. \smallskip

	La seguente tabella raccoglie le proprietà
	delle estensioni: \smallskip

	\begin{center}
		\scriptsize
		\vskip -0.1in
		\begin{tabular}{l|l|l|l|l}
			\hline
			$\mathcal{P}$ & Torri   & \textit{Shift} & Composto & Intersez. \\ \hline
			Est. fin.     & Strett. & Normal.        & Complet. & Sì        \\ \hline
			Est. alg.     & Strett. & Complet.       & Complet. & Sì        \\ \hline
			Est. sep.     & Strett. & Normal.        & Normal.  & Sì        \\ \hline
			Est. nor.     & Debolm. & Normal.        & Normal.  & Sì        \\ \hline
			Est. Gal.     & Debolm. & Normal.        & Normal.  & Sì
		\end{tabular}
	\end{center}

	\section{Teorema dell'elemento primitivo}

	Se $L / K$ è un'estensione finita e separabile,
	$L$ è in particolare un'estensione semplice di
	$K$, per il \textbf{Teorema dell'elemento primitivo}.
	In campi finiti, un tale elemento primitivo è
	un generatore di $L^*$. In campi infiniti, per
	$L = K(a, b)$,
	si può invece considerare il seguente polinomio:
	\[ p(x) = \prod_{i < j} (\varphi_i(a) + x \varphi_i(b) - \varphi_j(b) - x \varphi_j(b)), \]
	dove le varie $\varphi_i$ sono le $K$-immersioni di
	$L$ su $\overline{K}$.
	Si verifica che $p(x)$ è non nullo, e pertanto
	ha supporto non vuoto. Pertanto esiste un $t \in K$ tale
	per cui $p(t) \neq 0$, da cui si ricava che
	$L = K(a + bt)$ (perché ha il numero giusto di coniugati per
	generare tutto $L$!). Reiterando questo algoritmo su
	tutti i generatori dell'estensione, si ottiene
	un elemento primitivo desiderato.

	\section{Corrispondenza di Galois}

	\subsection{Teorema di corrispondenza di Galois}

	Se $L / K$ è di Galois, detto $H \leq \Gal(L / K)$,
	si definisce $L^H$ come la sottoestensione di $L$ composta
	dagli elementi di $L$
	fissati da tutti gli elementi di $H$.
	In particolare vale che $L^H = K \iff H = \Gal(L / K)$.
	Conseguentemente, vale il seguente teorema:

	\begin{theorem}[di corrispondenza di Galois]
		Sia $\mathcal{E}$ l'insieme delle sottoestensioni
		di $L / K$ estensione di Galois. Sia
		$\mathcal{G}$ l'insieme dei sottogruppi di
		$\Gal(L / K)$. Allora $\mathcal{E}$ è
		in bigezione con $\mathcal{G}$ attraverso
		la mappa $\alpha : \mathcal{E} \to \mathcal{G}$
		tale per cui:
		\[ F \xmapsto{\alpha} \Gal(L / F) \leq
			\Gal(L / K), \]
		la cui inversa $\beta : \mathcal{G} \to \mathcal{E}$
		è tale per cui:
		\[ H \xmapsto{\beta} L^H \subseteq L. \]
		Questa corrispondenza manda sottoestensioni di grado $n$ su $K$
		in sottogruppi di indice $n$ e viceversa (infatti $[F : K] = [L : K] / [L : F] = \abs{\Gal(L / K)} / \abs{\Gal(L : F)} =
			[\Gal(L / K) : \Gal(L / F)]$). \smallskip

		Inoltre, una sottoestensione $F / K$ di
		$L / K$ è normale su $K$ se e solo se
		il corrispondente sottogruppo di $\Gal(L / K)$
		è normale. Infine, se $F / K$ è normale,
		$F$ è in particolare di Galois e vale:
		\[ \Gal(F / K) \cong \faktor{\Gal(L / K)}{\Gal(L / F)}. \]

		Infine valgono le seguente proprietà:
		\begin{itemize}
			\item $L^H \subseteq L^Q \iff H \geq Q$ (la corrispondenza inverte le inclusioni),
			\item $L^H L^Q = L^H(L^Q) = L^{H \cap Q}$,
			\item $L^{\gen{H, Q}} = L^H \cap L^Q$,
		\end{itemize}
	\end{theorem}

	In particolare, un diagramma di campi -- a patto
	che il suo estremo superiore sia di Galois -- può
	essere collegato ad un diagramma di gruppi,
	``invertendo'' le inclusioni. Se
	$G = \Gal(L / K)$ e $H \subseteq G$, allora il
	diagramma:
	\[\begin{tikzcd}
			L \\
			{L^H} \\
			K
			\arrow["G", bend left, no head, from=3-1, to=1-1]
			\arrow["{G/H}"', no head, from=3-1, to=2-1]
			\arrow["H"', no head, from=2-1, to=1-1]
		\end{tikzcd}\]
	si relaziona tramite corrispondenza al
	diagramma:
	\[\begin{tikzcd}[row sep=small]
			{\{e\}} \\
			\\
			H \\
			\\
			G
			\arrow[no head, from=1-1, to=3-1]
			\arrow[no head, from=3-1, to=5-1]
		\end{tikzcd}\]

	\subsection{Relazioni tra gruppi di Galois in un diagramma}

	Consideriamo il seguente diagramma di campi:
	\[\begin{tikzcd}
			& LM \\
			L & {} & M \\
			& {L \cap M} \\
			& K
			\arrow[no head, from=4-2, to=3-2]
			\arrow[no head, from=3-2, to=2-1]
			\arrow[no head, from=2-1, to=1-2]
			\arrow[no head, from=3-2, to=2-3]
			\arrow[no head, from=2-3, to=1-2]
		\end{tikzcd}\]

	Allora $\Gal(LM/M)$ è sempre isomorfo a $\Gal(L/(L \cap M))$ tramite:
	\[
		\varphi \mapsto \restr{\varphi}{L}.
	\]
	Per la validità dell'isomorfismo, è sufficiente che $L$ sia di Galois, ma non
	è necessario che lo sia anche $M$ (benché debba essere finita). \smallskip

	Pertanto, se $L \cap M = K$, vale $\Gal(LM/M) \cong \Gal(L/K)$, e quindi
	in tal caso:
	\[
		[LM : K] = [L : K] [M : K].
	\]

	Sussiste sempre un'immersione:
	\[
		\Gal(LM/K) \hookrightarrow \Gal(L/K) \times \Gal(M/K),
	\]
	che manda $\varphi$ in $(\restr{\varphi}{L}, \restr{\varphi}{M})$. Tale
	immersione è un isomorfismo se e solo se $L \cap M = K$, per il precedente
	risultato.

	\subsection{Sottoestensioni quadratiche di un campo}

	Per il Teorema di corrispondenza di Galois, le sottoestensioni quadratiche di un'estensione
	di Galois $L/K$ sono in corrispondenza con i sottogruppi di indice $2$ di
	$\Gal(L/K)$. \smallskip

	Si osserva che un sottogruppo $H \leq G$ di indice $2$ contiene
	sempre il sottogruppo dei suoi quadrati $G^2 = \langle g^2 \mid g \in G \rangle$.
	Pertanto, i sottogruppi di indice $2$ di $G = \Gal(L/K)$ sono in corrispondenza a
	loro volta con i sottogruppi di indice $2$ di $G / G^2$, che è della forma:
	\[
		G/G^2 \cong (\ZZ / 2 \ZZ)^k,
	\]
	dal momento che ogni suo elemento ha ordine al più $2$. In $(\ZZ / 2 \ZZ)^k$
	vi sono esattamente $2^k-1$ iperpiani (i.e., sottogruppi di indice $2$), in
	quanto corrispondenti ai $2^k-1$ vettori non nulli dello spazio vettoriale. Pertanto
	$2^k-1$ è il numero di sottoestensioni quadratiche di $L$.

	\subsection{Realizzazione di un gruppo come gruppo di Galois}

	Se $p$ è un primo e $f(x) \in \QQ[x]$ è un irriducibile di grado $p$ con esattamente
	$2$ radici complesse non reali, allora $\Gal(L/\QQ) \cong S_p$, dove $L$ è il
	campo di spezzamento di $f(x)$ su $\QQ$. Infatti, esisterebbe
	in $\Gal(L/\QQ)$ una trasposizione indotta dal coniugio e un $p$-ciclo dal momento che
	$p \mid [L : \QQ]$; su $S_p$, una trasposizione qualsiasi e un $p$-ciclo generano tutto
	il gruppo, da cui il risultato desiderato. \smallskip

	Per il Lemma di Artin, se $K$ è un campo qualsiasi, e $G$ è un sottogruppo finito di $\Aut(K)$, allora
	$K/K^G$ e un'estensione di Galois finita con $\Gal(K/K^G) = G$. \smallskip

	Pertanto ogni gruppo finito $G$ si realizza come gruppo di Galois di un'estensione di campi. Infatti,
	se $|G| = n$, $G$ si immerge in $S_n$, e $S_n$ si immerge in $S_p$ per un primo $p \geq n$. Allora è
	sufficiente prendere un polinomio irriducibile di grado $p$ con esattamente $2$ radici
	complesse non reali su $\QQ[x]$ e considerare l'opportuna sottoestensione del suo campo di spezzamento
	tramite la corrispondenza di Galois.

	\subsection{Discriminante, gruppo alterno e gruppo di Galois di un irriducibile di grado \texorpdfstring{$3$}{3}}

	Dato un polinomio $f(x) \in K[x]$ di grado $n$ e le sue radici
	$a_1$, ..., $a_n$ in $\overline{K}$, si dice \textbf{discriminante}
	$\Delta(f)$ la seguente quantità:
	\[
		\Delta(f) = \prod_{i < j} (a_i - a_j)^2.
	\]
	Il polinomio $f(x)$ ha radici multiple se e solo se
	$\Delta(f)$ si annulla. Inoltre, $\Delta(f)$ è un polinomio
	simmetrico nelle radici, e quindi, per il Teorema fondamentale
	dei polinomi simmetrici, si può scrivere nella base degli
	$e_i(a_1, \ldots, a_n)$, che sono, a meno di segno e fattore di
	riscalamento, i coefficienti del polinomio stesso (e quindi
	$\Delta(f) \in K$). \smallskip

	Il discriminante è invariante per traslazione. I primi discriminanti
	sono calcolati come segue:
	\begin{itemize}
		\item il polinomio $x^2 + bx + c$ ha discriminante $b^2 - 4c$,
		\item il polinomio $x^3 + px + q$ ha discriminante $-4p^3-27q^2$.
	\end{itemize}

	Supponiamo ora $f(x)$ irriducibile. Sia $L$ il campo di spezzamento di
	$f(x)$ su $K$.
	Essendo una funzione simmetrica nelle radici, il discriminante è invariante
	rispetto all'azione del gruppo $\Gal(L/K)$, (questo mostra ancora che
	$\Delta(f) \in K$ per corrispondenza). \smallskip

	Si osserva che, per $\sigma \in \Gal(L/K)$:
	\[
		\sigma\left(\prod_{i < j} (a_i - a_j)\right) = \prod_{i < j} (a_{\sigma(i)} - a_{\sigma(j)}) = \sgn(\sigma) \prod_{i < j} (a_i - a_j),
	\]
	dove si è trasferito l'azione sugli indici delle radici. Pertanto,
	$\Gal(L/K)$ si immerge nel gruppo alterno $A_n$ se e solo se
	la radice quadrata di $\Delta(f)$ è fissata, ovverosia se e solo se
	$\Delta(f)$ è un quadrato in $K$. \smallskip

	Se $n = 3$, allora, dal momento che le uniche possibilità per cardinalità sono
	$A_3 \cong \ZZ / 3 \ZZ$ e $S_3$, vale che:
	\[
		\Gal(L/K) \cong \begin{cases}
			\ZZ / 3 \ZZ & \text{se } \Delta(f) = -4p^3-27q^2 \text{ quadrato in } K, \\
			S_3         & \text{altrimenti}.
		\end{cases}
	\]

	\section{Gruppi di Galois noti}

	\subsection{Campi finiti}

	Il campo finito $\FF_{p^n}$ è sempre normale
	su $\FF_p$, dal momento che può essere costruito
	come campo di spezzamento di $x^{p^n} - x$ su
	$\FF_p$ stesso. Equivalentemente, poiché
	un omomorfismo di campi è sempre iniettivo (e dunque
	conserva sempre la cardinalità),
	una $\FF_p$-immersione deve mandare $\FF_{p^n}$
	in un campo della stessa cardinalità, e quindi
	necessariamente un campo isomorfo a $\FF_{p^n}$. \smallskip

	Per un campo finito, $\Frob$ è un automorfismo che
	fissa $\FF_p$. Allora $\Frob \in \Gal(\FF_{p^n} / \FF_p)$.
	Inoltre $\ord \Frob = n = \abs{\Gal(\FF_{p^n} / \FF_p)}$ (altrimenti $\FF_{p^n}$ non sarebbe campo di
	spezzamento di $x^{p^n}-x$), e quindi vale che:
	\[ \Gal(\FF_{p^n} / \FF_p) = \gen{\Frob} \cong \ZZmod{n}. \]

	Pertanto se $\alpha \in \FF_{p^n} \setminus \FF_p$,
	tutti i suoi coniugati si ottengono reiterando
	al più $p^n$ volte $\Frob$ su $\alpha$. \smallskip

	Se il campo base non è $\FF_p$, si può utilizzare la corrispondenza
	di Galois per determinare comunque che:
	\[
		\Gal(\FF_{p^n} / \FF_{p^m}) \cong \faktor{\ZZ / n \ZZ}{\ZZ / m \ZZ}.
	\]

	\subsection{Polinomi biquadratici}

	Sia $p(x) = x^4 + ax^2 + b$ irriducibile su $\QQ$.
	Allora, se $L$ è un suo campo di spezzamento e $\Delta = a^2 - 4b$ è l'usuale discriminante di $p$ visto come polinomio in $x^2$, vale che:

	\[ \Gal(L / \QQ) \cong \begin{cases}
			\ZZmod{4}                  & \se b \Delta \text{ è quadrato in $\QQ$}, \\
			\ZZmod{2} \times \ZZmod{2} & \se b \text{ è quadrato in $\QQ$},        \\
			D_4                        & \altrimenti.
		\end{cases} \]

	\subsection{Radici di primi in \texorpdfstring{$\QQ$}{ℚ}}

	Siano $p_1$, ..., $p_n$ numeri primi distinti.
	Allora vale che:
	\[ \Gal(\QQ(\sqrt{p_1}, \ldots, \sqrt{p_n})/\QQ) \cong (\ZZmod{2})^n. \]
	Si possono infatti contare le sottoestensioni quadratiche di
	$\QQ(\sqrt{p_1}, \ldots, \sqrt{p_{n-1}})$ e constatare che sono esattamente
	le $\QQ(\sqrt{p_1^{\eps_1} \cdots p_{n-1}^{\eps_{n-1}}})$ con $\eps_i \in \{\pm 1\}$
	e almeno un $\eps_i$ non nullo. Dal momento che nessuna di questa può coincidere con
	$\QQ(\sqrt{p_n})$, si ottiene il risultato desiderato. \smallskip

	Un generatore di questa estensione è $\sqrt{p_1} + \ldots + \sqrt{p_n}$ dacché
	ha $2^n$ coniugati.

	\subsection{Radici dell'unità e polinomi ciclotomici \texorpdfstring{$\Phi_n(x)$}{Φₙ(x)}}

	Dato un campo $K$, il polinomio $x^n - 1$ è separabile se e solo se
	$\Char K \nmid n$. Infatti, per il Criterio della derivata, $x^n-1$
	è separabile se e solo se $\MCD(nx^{n-1}, x^n - 1) = 1$.

	Sia $\Phi_n(x)$ l'$n$-esimo polinomio ciclotomico, così definito:
	\[ \Phi_n(x) = \prod_{\substack{1 \leq d \leq n \\ \MCD(d, n) = 1}} (x - \zeta_n^d), \]
	dove $\zeta_n$ è una radice primitiva $n$-esima dell'unità. \smallskip

	Di seguito sono presentati i primi $8$ polinomi ciclotomici:
	\begin{itemize}
		\item $\Phi_1(x) = x - 1$,
		\item $\Phi_2(x) = x + 1$,
		\item $\Phi_3(x) = x^2 + x + 1$,
		\item $\Phi_4(x) = x^2 + 1$,
		\item $\Phi_5(x) = x^4 + x^3 + x^2 + x + 1$,
		\item $\Phi_6(x) = x^2 - x + 1$,
		\item $\Phi_7(x) = x^6 + x^5 + x^4 + x^3 + x^2 + x + 1$,
		\item $\Phi_8(x) = x^4 + 1$.
	\end{itemize}

	Tale polinomio è sempre a coefficienti interi ed è inoltre primitivo
	su $\ZZ[x]$. Vale inoltre che:
	\[ x^n - 1 = \prod_{m \mid n} \Phi_m(x). \]
	Il campo di spezzamento di $\Phi_n(x)$ su $\QQ$ è
	$\QQ(\zeta_n)$, che è quindi un'estensione normale e separabile,
	e pertanto di Galois finita. \smallskip

	Inoltre vale che:
	\[ \Gal(\QQ(\zeta_n)/\QQ) \cong (\ZZmod{n})^\times, \]
	e $\Phi_n(x)$ è sempre irriducibile su $\QQ$. \smallskip

	In generale, vale per $n \geq 3$:
	\[
		[\QQ(\zeta_n) : \QQ(\zeta_n) \cap \RR] = 2,
	\]
	ovverosia il grado di $\QQ(\zeta_n)$ sulla sua parte reale è $2$.
	Infatti, $\alpha = \zeta_n + \zeta_n\inv \in \QQ(\zeta_n) \cap \RR$ (dacché
	$\overline{\zeta_n} = \zeta_n\inv$) e $\zeta_n$ è radice di
	$x^2 - \alpha x + 1$. Si verifica inoltre che
	$\QQ(\zeta_n) \cap \RR$ coincide proprio con $\QQ(\zeta_n + \zeta_n\inv)$, che è
	estensione normale di $\QQ$.

	\vfill
	\hrule
	~\\
	Ad opera di Gabriel Antonio Videtta, \url{https://poisson.phc.dm.unipi.it/~videtta/}.
	~\\Reperibile su
	\url{https://github.com/hearot/notes}, nella sezione \textit{Corsi $\to$ Algebra 1 $\to$ 3. Teoria delle estensioni di campo e di Galois $\to$ Scheda riassuntiva di Teoria dei campi e di Galois}.
\end{multicols}

\end{document}